\chapter{P\&L et comptabilité d'une position}\label{ch:m1:pnl-and-positions}

À 16:05, la trader dit à son responsable de desk qu'elle a gagné quarante mille dollars. À
17:30, le rapport de risque indique trente et un. Le lendemain matin, la direction
financière comptabilise vingt-sept. Personne ne ment et personne ne s'est trompé~:
les trois chiffres s'appuient sur les mêmes six transactions et donnent trois réponses
différentes aux questions «~à quel prix la \omterm{def:m1:pnl-and-positions:position}{position} vaut-elle~?~» et «~combien la journée
a-t-elle coûté~?~». Une société de trading vit et rémunère ses équipes sur ce chiffre. Ce
chapitre le définit, montre quelle partie relève de l'arithmétique et quelle partie
relève de la convention, et construit le composant qui le calcule.

\section{Positions et exposition}

\begin{definition}[Position]\label{def:m1:pnl-and-positions:position}
Une \emph{position}\index{position} sur un instrument est la quantité algébrique
détenue~: positive (\emph{long}) si l'on en est propriétaire, négative (\emph{short}) si l'on en est redevable. Elle
ne change que par exécutions, opérations sur titres et transferts~; elle est la somme
des quantités algébriques de tous ces mouvements depuis l'origine.
\end{definition}

\begin{definition}[Notionnel, exposition brute et exposition nette]\label{def:m1:pnl-and-positions:exposure}
Le \emph{notionnel}\index{notionnel} d'une \omterm{def:m1:pnl-and-positions:position}{position} $q_i$ valorisée au prix $p_i$
est $|q_i|\,p_i$ (multiplié par le multiplicateur du contrat, pour un dérivé). Pour un
portefeuille, l'\emph{exposition brute}\index{exposition brute} est $\sum_i |q_i| p_i$
et l'\emph{exposition nette}\index{exposition nette} est $\sum_i q_i p_i$.
\end{definition}

\begin{example}[Un portefeuille long--short]\label{ex:m1:pnl-and-positions:book}
Long de 300~millions de dollars et short de 200~millions de dollars, pour 100~millions de dollars de fonds propres~:
\omterm{def:m1:pnl-and-positions:exposure}{exposition brute} de 500~millions de dollars, soit $5\times$ les fonds propres, mesure de l'ampleur des pertes possibles~;
\omterm{def:m1:pnl-and-positions:exposure}{exposition nette} de 100~millions de dollars, soit $1\times$, mesure de la direction de marché
effectivement prise. Des limites sont fixées sur les deux mesures, et sur chaque \omterm{def:m1:pnl-and-positions:position}{position} séparément.
\end{example}

\section{Valorisation au prix de marché}

\begin{definition}[Valorisation au prix de marché]\label{def:m1:pnl-and-positions:mtm}
La \emph{valorisation au prix de marché}\index{valorisation au prix de marché} (mark to market) d'une \omterm{def:m1:pnl-and-positions:position}{position} consiste à la valoriser à un prix de marché
courant, le \emph{mark}, plutôt qu'au montant payé pour elle.
Le P\&L d'un portefeuille sur une période, c'est-à-dire son profit ou sa perte, est la variation
de sa valeur au prix de marché, augmentée de la trésorerie nette qu'il a produite.
\end{definition}

\begin{proposition}[L'identité du P\&L]\label{prop:m1:pnl-and-positions:identity}
On part d'une \omterm{def:m1:pnl-and-positions:position}{position} nulle. Après les exécutions $k = 1,\dots,n$, de sens $\epsilon_k = \pm1$ (achat
$+1$), quantités $n_k$, prix $p_k$ et frais $f_k$, la \omterm{def:m1:pnl-and-positions:position}{position} est $q =
\sum_k \epsilon_k n_k$, la trésorerie est $c = -\sum_k \epsilon_k n_k p_k$, et le
P\&L au mark $m$ est
\[
\pnl(m) \;=\; c + q\,m - \sum_k f_k .
\]
Il ne dépend des exécutions qu'à travers $(q, c, \sum f_k)$~: d'aucune convention
comptable, et pas non plus de l'ordre des exécutions.
\end{proposition}
\begin{proof}
Le portefeuille détient la trésorerie $c - \sum f_k$ et $q$ unités valant chacune $m$~; il
était parti de zéro.
\end{proof}

C'est la partie exacte. Tout est entier si les prix sont exprimés en ticks ou en unités
comptables (voir la \cref{bld:m1:what-a-trading-firm-does:ledger}), et un système de P\&L
qui ne la reproduit pas à l'unité près a perdu une transaction. La
\emph{ventilation} de ce chiffre est l'endroit où entrent les conventions.

\begin{definition}[Coût moyen, P\&L réalisé et P\&L latent]\label{def:m1:pnl-and-positions:realised}
Selon la convention du \emph{coût moyen}\index{coût moyen}, une \omterm{def:m1:pnl-and-positions:position}{position} ouverte
porte le prix moyen pondéré par les quantités $\bar p$ des exécutions qui l'ont constituée.
Une exécution qui réduit la \omterm{def:m1:pnl-and-positions:position}{position} de $n$ unités au prix $p$
\emph{réalise} $\pm n(p - \bar p)$ (plus pour une \omterm{def:m1:pnl-and-positions:position}{position} longue, moins pour une \omterm{def:m1:pnl-and-positions:position}{position} courte) et
laisse $\bar p$ inchangé~; une exécution qui l'augmente met $\bar p$ à jour~; une exécution
qui inverse la \omterm{def:m1:pnl-and-positions:position}{position} clôture entièrement l'ancienne et ouvre le reliquat au
prix $p$. Le \emph{P\&L réalisé}\index{P\&L réalisé} est la somme de ces
montants~; le \emph{P\&L latent}\index{P\&L latent} est $q(m - \bar
p)$.
\end{definition}

\begin{example}[Six exécutions]\label{ex:m1:pnl-and-positions:six}
La trader du paragraphe d'ouverture achète 20\,000 à 49.30 et 20\,000 à
49.55 ($\bar p = 49.425$), vend 10\,000 à 50.10 (réalisant $10\,000 \times
0.675 = 6\,750$), achète 10\,000 à 49.70 ($\bar p = 49.494$), vend 25\,000 à
50.28 (réalisant $19\,656$), et achète 15\,000 à 50.20 ($\bar p = 49.847$).
Elle termine longue de 30\,000 avec 26\,406~\$ réalisés. Selon la méthode \emph{premier entré, premier
sorti}, les mêmes ventes sont appariées aux achats les plus anciens et réalisent
28\,750~\$. Valorisées à 50.00, les parts latentes sont de 4\,594~\$ et 2\,250~\$~:
les deux conventions totalisent 31\,000~\$, comme l'exige la \cref{prop:m1:pnl-and-positions:identity}.
La ventilation compte pour l'impôt et pour les statistiques de performance
(«~taux de réussite~»)~; elle ne change jamais la valeur de la société.
\end{example}

\begin{omfigure}
\begin{tikzpicture}
\begin{axis}[width=11cm,height=5.2cm,
  xlabel={},ylabel={milliers de dollars},
  xtick={0,1,2,3,4,5,6,7},xticklabels={09:30,09:35,10:10,11:20,13:05,14:30,15:40,16:00},
  tick label style={font=\footnotesize},label style={font=\small},
  xmin=0,xmax=7,ymin=-2,ymax=42,grid=major,grid style={gray!25},
  legend columns=3,legend style={font=\footnotesize,at={(0.5,-0.22)},anchor=north,draw=none,fill=none,/tikz/every even column/.append style={column sep=8pt}},
  table/col sep=comma]
\addplot[omMeth,very thick,const plot] table[x=k,y=realised_k]{figdata/markets-1/07-pnl-and-positions/day.csv};
\addplot[omProp,very thick,mark=*,mark size=1.2pt] table[x=k,y=unrealised_k]{figdata/markets-1/07-pnl-and-positions/day.csv};
\addplot[omDef,very thick,mark=*,mark size=1.4pt] table[x=k,y=total_k]{figdata/markets-1/07-pnl-and-positions/day.csv};
\legend{réalisé (coût moyen),latent,total au mid}
\end{axis}
\end{tikzpicture}

\omcaption{La journée de l'\cref{ex:m1:pnl-and-positions:six}, observée à chaque
exécution et à la clôture, valorisée au mid. Le \omterm{def:m1:pnl-and-positions:realised}{P\&L réalisé} ne varie que lorsqu'une
vente réduit la \omterm{def:m1:pnl-and-positions:position}{position}~; le dernier point baisse parce que le mid de clôture est
inférieur aux cours de l'après-midi alors qu'elle est toujours longue de 30\,000 actions. Données~:
la liste des exécutions du chapitre.}\label{fig:m1:pnl-and-positions:day}
\end{omfigure}

\begin{method}[Choisir une valorisation]\label{met:m1:pnl-and-positions:mark}
\begin{enumerate}
\item \textbf{La dernière transaction} est l'information la plus récente et la moins bonne
  base de valorisation~: elle se situe au cours acheteur ou au cours vendeur, pour une taille qui n'est pas la sienne, éventuellement
  datant de quelques minutes.
\item \textbf{Le mid} d'une cotation serrée et bilatérale est la valorisation intrajournalière usuelle
  pour les instruments liquides.
\item \textbf{Le côté défavorable} --- cours acheteur pour les \omterm{def:m1:pnl-and-positions:position}{positions} longues, cours vendeur pour les \omterm{def:m1:pnl-and-positions:position}{positions} courtes ---
  valorise la \omterm{def:m1:pnl-and-positions:position}{position} au prix que rapporterait sa clôture pour une petite taille~; pour une
  \omterm{def:m1:pnl-and-positions:position}{position} importante, il faut retrancher en outre une estimation de l'impact.
\item \textbf{La clôture officielle} (prix de l'enchère de clôture ou prix de règlement-livraison)
  est celle qu'utilisent la finance, les chambres de compensation et les administrateurs de fonds~: elle est
  publique, unique et auditable.
\end{enumerate}
Décider une fois pour toutes, par écrit, quel cours sert à quel rapport. Un trader autorisé à
choisir la valorisation peut choisir le P\&L.
\end{method}

\begin{example}[Quarante, trente et un, vingt-sept]\label{ex:m1:pnl-and-positions:three}
L'écran du trader valorise au cours de la dernière transaction, 50.30~: la trésorerie $-\$1\,469\,000$
plus $30\,000 \times 50.30$ donne 40\,000~\$. Le système de risque valorise au
mid de clôture, 50.00~: 31\,000~\$. L'écart de 9\,000~\$ entre les deux correspond à $30\,000$ actions
$\times$ 30 cents de valorisation. La finance retient le cours de clôture officiel, également 50.00, et
soustrait 3\,400~\$ de frais et 600~\$ de financement~: 27\,000~\$
(voir la \cref{fig:m1:pnl-and-positions:three}).
\end{example}

\begin{omfigure}
\begin{tikzpicture}
\begin{axis}[width=11cm,height=3.6cm,xbar,bar width=8pt,
  xlabel={P\&L des mêmes six exécutions (milliers de dollars)},
  ytick=data,yticklabels={{Trader (dernier cours~; brut)},{Risques (prix milieu de clôture~; brut)},{Finance (clôture officielle~; net)}},
  yticklabel style={font=\small},tick label style={font=\small},label style={font=\small},
  y dir=reverse,xmin=0,xmax=45,enlarge y limits=0.25,grid=major,grid style={gray!25},
  nodes near coords,nodes near coords style={font=\footnotesize},
  table/col sep=comma]
\addplot[fill=omDef!60,draw=omDef] table[x=usd_k,y=k]{figdata/markets-1/07-pnl-and-positions/three_numbers.csv};
\end{axis}
\end{tikzpicture}

\omcaption{Une journée, trois nombres. Données~: l'\cref{ex:m1:pnl-and-positions:three},
calculées par le script du chapitre.}\label{fig:m1:pnl-and-positions:three}
\end{omfigure}

\section{Ce qu'une position coûte et rapporte pendant sa détention}

\begin{definition}[Carry]\label{def:m1:pnl-and-positions:carry}
Le \emph{carry}\index{carry} d'une \omterm{def:m1:pnl-and-positions:position}{position} est le P\&L qu'elle produit si les cours
ne bougent pas~: les revenus perçus (dividendes, coupons, intérêts sur le produit des ventes à découvert)
moins les coûts payés (financement des \omterm{def:m1:pnl-and-positions:position}{positions} longues, \omterm{def:m1:financing:seclending}{frais d'emprunt} des \omterm{def:m1:pnl-and-positions:position}{positions} courtes, dividendes dus sur les
\omterm{def:m1:pnl-and-positions:position}{positions} courtes).
\end{definition}

Un P\&L quotidien complet compte donc cinq postes, en dehors des mouvements de cours~:
les commissions et frais de bourse sur les exécutions du jour~; le financement couru
(voir la \cref{bld:m1:financing:calc})~; \omterm{def:m1:financing:seclending}{frais d'emprunt}~; dividendes et coupons, comptabilisés
à la date de détachement (voir au \cref{ch:m1:shares-corporate-actions-indices})~; et, pour
les \omterm{def:m1:pnl-and-positions:position}{positions} libellées en devise étrangère, la variation du taux de change.

\begin{proposition}[Conversion de devise]\label{prop:m1:pnl-and-positions:fx}
Une \omterm{def:m1:pnl-and-positions:position}{position} de $q$ unités d'un actif étranger passe du cours $p_0$ à $p_1$
dans sa propre devise, tandis que le taux de change (devise de base par unité de
devise étrangère) passe de $x_0$ à $x_1$. Son P\&L en devise de base est
\[
q\,(p_1x_1 - p_0x_0) \;=\; \underbrace{q\,(p_1-p_0)\,x_0}_{\text{cours}}
 \;+\; \underbrace{q\,p_0\,(x_1-x_0)}_{\text{devise}}
 \;+\; \underbrace{q\,(p_1-p_0)(x_1-x_0)}_{\text{croisé}} .
\]
\end{proposition}
\begin{proof}
Développer $p_1x_1 = (p_0 + \Delta p)(x_0 + \Delta x)$.
\end{proof}

\begin{example}[Une bonne action dans une mauvaise devise]\label{ex:m1:pnl-and-positions:fx}
Un fonds en dollars détient 10\,000 actions d'une valeur en euros qui passe de 40
à 41, tandis que l'euro passe de 1.10 à 1.08 dollars. Effet cours~: $+\$11\,000$.
Effet de change~: $-\$8\,000$. Effet croisé~: $-\$200$. Total $+\$2\,800$. Un fonds qui veut
le premier chiffre et non le second vend 400\,000~\euro à terme~; la couverture
couvre la valeur d'ouverture $qp_0$, jamais le terme croisé, et doit être redimensionnée à mesure
que l'action évolue.
\end{example}

\section{P\&L explain~: une première approche}

\begin{definition}[Attribution du P\&L]\label{def:m1:pnl-and-positions:attribution}
Une \emph{attribution du P\&L}\index{attribution du P\&L}, ou \emph{explain}, est une
décomposition du P\&L d'une période en causes nommées dont la somme égale le
total, tout reliquat étant signalé comme \emph{non expliqué}.
\end{definition}

\begin{method}[L'explain quotidien d'un livre de trésorerie]\label{met:m1:pnl-and-positions:explain}
Avec $q_0$ la \omterm{def:m1:pnl-and-positions:position}{position} à la clôture précédente $m_0$, et $m_1$ la clôture du jour~:
\begin{enumerate}
\item \textbf{\omterm{def:m1:pnl-and-positions:position}{Position} reportée~:} $q_0(m_1 - m_0)$.
\item \textbf{Nouvelles transactions~:} $\sum_k \epsilon_k n_k (m_1 - p_k)$ --- chaque exécution
  mesurée depuis son propre prix jusqu'à la clôture.
\item \textbf{Frais}, \textbf{financement}, \textbf{dividendes},
  \textbf{devise}~: chacun sur sa propre ligne.
\item \textbf{Non expliqué~:} total moins la somme des postes ci-dessus. Pour un livre de trésorerie
  il doit être nul au centime près~; tout autre écart signale une transaction manquante, une valorisation
  fausse ou une opération sur titre non traitée.
\end{enumerate}
\end{method}

L'identité qui sous-tend la méthode est
la \cref{prop:m1:pnl-and-positions:identity} appliquée entre deux clôtures~: les
deux premiers postes donnent $c + q_1 m_1 - q_0 m_0$. Le second poste est celui où se voit le savoir-faire
d'un desk de trading~: le P\&L des nouvelles transactions d'un \omterm{def:m1:what-a-trading-firm-does:market-maker}{teneur de marché} est sa capture de fourchette et
sa \omterm{def:m1:what-a-trading-firm-does:adverse-selection}{sélection adverse} (voir la \cref{prop:m1:what-a-trading-firm-does:capture})~; son
P\&L de \omterm{def:m1:pnl-and-positions:position}{position} reportée doit être faible et relève du pur risque. Pour les dérivés, le premier poste
est lui-même décomposé par facteur de risque, ce qui fait l'objet du
livre~6 de la série One Quant.

\begin{omfigure}
\begin{tikzpicture}
\begin{axis}[width=11cm,height=4.4cm,xbar,bar width=8pt,
  xlabel={milliers de \$},
  ytick=data,yticklabels={{position reportée},{nouvelles transactions},{frais},{financement},{total}},
  yticklabel style={font=\small},tick label style={font=\small},label style={font=\small},
  y dir=reverse,xmin=-4,xmax=14,enlarge y limits=0.15,grid=major,grid style={gray!25},
  nodes near coords,nodes near coords style={font=\footnotesize},
  table/col sep=comma]
\addplot[fill=omDef!60,draw=omDef] table[x=usd_k,y=k]{figdata/markets-1/07-pnl-and-positions/explain.csv};
\end{axis}
\end{tikzpicture}

\omcaption{L'explain du lendemain~: 30\,000 actions reportées de 50.00 à 50.40,
10\,000 vendues à 50.35 et 5\,000 achetées à 50.10, 900~\$ de frais et 650~\$ de
financement. La majeure partie des 11\,450~\$ de la journée a été dégagée par la \omterm{def:m1:pnl-and-positions:position}{position} de la veille, et non
par les décisions du jour. Données~: calculées par le script du chapitre.}
\end{omfigure}

\begin{omfigure}
\begin{tikzpicture}[font=\small,
  box/.style={draw=omDef,rounded corners=2pt,minimum width=2.6cm,minimum height=0.9cm,align=center,fill=omDef!4},
  src/.style={draw=omExo,rounded corners=2pt,minimum width=2.6cm,minimum height=0.9cm,align=center,fill=omExo!6},
  flow/.style={-{Stealth[length=2mm]},thick,omDef}]
\node[src] (fills) at (0,1.2) {Exécutions\\(drop copy)};
\node[src] (marks) at (0,-1.2) {Valorisations\\(politique par relevé)};
\node[box] (keeper) at (4.6,0) {Gestionnaire des positions\\et du P\&L};
\node[box] (explain) at (9.2,1.2) {Explain};
\node[box] (ledger) at (9.2,-1.2) {Grand livre};
\node[src] (recon) at (13.4,0) {Rapprochement\\avec le courtier};
\draw[flow] (fills) -- (keeper);
\draw[flow] (marks) -- (keeper);
\draw[flow] (keeper) -- (explain);
\draw[flow] (keeper) -- (ledger);
\draw[flow] (ledger) -- (recon);
\draw[flow] (explain) -- (recon);
\end{tikzpicture}

\omcaption{Où se place le gestionnaire dans la petite société. Les exécutions arrivent du
drop copy de la plateforme ou du courtier, jamais de la croyance de la stratégie sur
ce qu'elle a négocié~; les \omterm{def:m1:pnl-and-positions:position}{positions} et les liquidités sont rapprochées chaque matin
avec le relevé du courtier compensateur.}
\end{omfigure}

\section{Tutoriel~: trois nombres, une journée}\label{tut:m1:pnl-and-positions:three}

\begin{tutorial}
\textbf{Objectif.} Rejouer les six exécutions, reproduire 40\,000~\$, 31\,000~\$ et
27\,000~\$, et vérifier que la convention comptable n'a pas d'importance.
\textbf{État final~:} \cref{fig:m1:pnl-and-positions:day,fig:m1:pnl-and-positions:three}.
\begin{tutsteps}
\item \textbf{Le cœur du gestionnaire}, tel qu'il est écrit pour le côté systèmes de
la société. La quantité, les liquidités et les frais sont des entiers~; seule la
répartition publiée est en virgule flottante.
\omcode{code/firm/pnl/cpp/firm_pnl.hpp}{21}{45}{Tenue des positions au coût moyen en C++~: ajout, réduction, inversion~; et total exact.}\label{lst:m1:pnl-and-positions:cpp}
\item \textbf{Rejouer et valoriser de trois façons.}
\omcode{code/markets-1/07-pnl-and-positions/python/pnl_day.py}{28}{41}{Le même objet de position, trois valorisations et deux lignes de coût.}\label{lst:m1:pnl-and-positions:three}
\item \textbf{Lancer les tests}~: \texttt{make test-code CH=markets-1/07-pnl-and-positions}.
Ils vérifient les trois nombres, et que le premier entré, premier sorti donne le même
total avec une répartition différente (28\,750~\$ réalisés contre 26\,406~\$).
\item \textbf{Mettre à l'épreuve l'identité.} Les tests de la construction rejouent cinq mille
exécutions aléatoires et exigent que la répartition coïncide avec le total exact.
\end{tutsteps}
\textbf{Pour aller plus loin.} Valoriser la \omterm{def:m1:pnl-and-positions:position}{position} longue au cours acheteur, 49.99~: quel est le
quatrième nombre~? Ajouter ensuite un dividende de 20 cents détaché au cours de la journée et
décider à quelle ligne de l'explain il appartient.
\end{tutorial}

\section{Construction~: le gestionnaire de positions et de P\&L}\label{bld:m1:pnl-and-positions:keeper}

\begin{build}
\bfield{Rôle} Le composant que demandent toutes les stratégies, tous les contrôles de risque et tous les rapports de
la petite société~: que détenons-nous, et combien vaut-il~? C'est la première
pièce construite dans les trois langages, car elle s'exécute à la fois dans le
processus à faible latence (livre~13) et dans la
pile de recherche et de reporting.
\bfield{Interface} \texttt{Position.on\_fill(side, quantity, price, fee)},
\texttt{unrealised(mark)}, \texttt{total(mark)}; \texttt{Book.on\_fill(symbol,
\dots)}, \texttt{total(marks)}, \texttt{gross\_exposure(marks)},
\texttt{net\_exposure(marks)}. Les prix, les liquidités et les frais sont des entiers en unités
comptables.
\bfield{Règles} \texttt{total} est calculé à partir de la quantité, des liquidités et des frais seulement,
et il est exact. La répartition entre \omterm{def:m1:pnl-and-positions:realised}{P\&L réalisé} et \omterm{def:m1:pnl-and-positions:realised}{P\&L latent} suit
la \cref{def:m1:pnl-and-positions:realised}. Les quantités ou prix non positifs
et les frais négatifs sont rejetés. Aucune allocation et aucune exception sur le
chemin d'exécution en C++ et en Rust, sauf pour les entrées rejetées.
\bfield{Tests d'acceptation} \texttt{code/firm/pnl/tests/}, \texttt{cpp/} et
\texttt{rust/}~: aller-retour~; \omterm{def:m1:pnl-and-positions:realised}{coût moyen} et clôture partielle~; inversion~; cinq
mille exécutions aléatoires dont la répartition égale le total~; rejets. Les trois
implémentations doivent concorder sur le même fichier d'exécutions.
\bfield{Pour aller plus loin} Ajouter des lots premier entré, premier sorti derrière la même interface et un
commutateur par compte~; publier chaque jour les lignes de l'explain dans le grand livre du
\cref{ch:m1:what-a-trading-firm-does}.
\end{build}

\begin{omsources}
\item L.~Harris, \emph{Trading and Exchanges}, Oxford University Press, 2003,
chapitre 21 (évaluation de la performance).
\item R.~Grinold et R.~Kahn, \emph{Active Portfolio Management}, 2e éd.,
McGraw-Hill, 2000, chapitre 17 (analyse de la performance).
\item International Accounting Standards Board, \emph{IFRS 13 Fair Value
Measurement} --- la norme comptable qui fonde le choix des valorisations.
\end{omsources}

\section{Exercices}

\begin{exercise}[$\star$]\label{exo:m1:pnl-and-positions:1}
Un livre est long de 40\,000 actions à 25~\$, long de 10\,000 à 120~\$ et court de
30\,000 à 50~\$, pour des fonds propres de 2~millions de dollars. Calculer son \omterm{def:m1:pnl-and-positions:exposure}{exposition brute} et son \omterm{def:m1:pnl-and-positions:exposure}{exposition nette} en
dollars et en multiples des fonds propres.
\end{exercise}

\begin{exercise}[$\star$]\label{exo:m1:pnl-and-positions:2}
Partant de zéro~: acheter 500 à 20.00, acheter 300 à 20.40, vendre 600 à 20.50. Donner
le \omterm{def:m1:pnl-and-positions:realised}{coût moyen} après chaque exécution, le \omterm{def:m1:pnl-and-positions:realised}{P\&L réalisé} et le \omterm{def:m1:pnl-and-positions:realised}{P\&L latent}
à une valorisation de 20.30.
\end{exercise}

\begin{exercise}[$\star$]\label{exo:m1:pnl-and-positions:3}
Pour les exécutions de l'exercice précédent, calculer les flux de trésorerie, et vérifier
la \cref{prop:m1:pnl-and-positions:identity} à la valorisation de 20.30.
\end{exercise}

\begin{exercise}[$\star\star$]\label{exo:m1:pnl-and-positions:4}
Partant de zéro~: vendre 1\,000 à 80.00, vendre 1\,000 à 81.00, acheter 3\,000 à
79.50. Suivre la \omterm{def:m1:pnl-and-positions:position}{position}, le \omterm{def:m1:pnl-and-positions:realised}{coût moyen} et le \omterm{def:m1:pnl-and-positions:realised}{P\&L réalisé} à travers le
renversement, et donner le total à une valorisation de 79.80.
\end{exercise}

\begin{exercise}[$\star\star$]\label{exo:m1:pnl-and-positions:5}
Un fonds libellé en livres sterling détient 50\,000 actions d'un titre américain. Sur un mois, le cours
passe de 200~\$ à 190~\$ et le dollar de 0.80 à 0.84 livre.
Calculer les termes de prix, de devise et croisé, ainsi que le total en livres.
\end{exercise}

\begin{exercise}[$\star\star$]\label{exo:m1:pnl-and-positions:6}
Le cours de clôture de la veille était de 30.00, avec une \omterm{def:m1:pnl-and-positions:position}{position} de $-20\,000$ actions. Aujourd'hui le
desk achète 8\,000 à 29.60 et vend 5\,000 à 29.90~; la clôture est à 29.70~; les frais
s'élèvent à 260~\$, les \omterm{def:m1:financing:seclending}{frais d'emprunt} courus à 120~\$ et les produits de la vente à découvert
ont rapporté 65~\$ d'intérêts. Produire l'explication du P\&L et la \omterm{def:m1:pnl-and-positions:position}{position} de clôture.
\end{exercise}

\begin{exercise}[$\star\star\star$]\label{exo:m1:pnl-and-positions:7}
\emph{Programmation.} Implémenter le premier entré, premier sorti (FIFO) pour des \omterm{def:m1:pnl-and-positions:position}{positions} pouvant être
courtes comme longues (le \texttt{fifo} du chapitre ne traite qu'une journée en \omterm{def:m1:pnl-and-positions:position}{position} longue uniquement). Le tester
sur les exécutions de l'exercice~\ref{exo:m1:pnl-and-positions:4} et indiquer le
\omterm{def:m1:pnl-and-positions:realised}{P\&L réalisé}~; le comparer au \omterm{def:m1:pnl-and-positions:realised}{coût moyen}.
\end{exercise}

\begin{exercise}[$\star\star\star$]\label{exo:m1:pnl-and-positions:8}
\emph{Trouver le défaut.} Un desk annonce un taux de réussite de 78~\% «~sur les transactions clôturées~»
et un \omterm{def:m1:pnl-and-positions:realised}{P\&L réalisé} positif pendant onze mois consécutifs~; son P\&L total
sur la période est négatif. Expliquer comment les deux peuvent être vrais, quelle convention
le rend facile, et quel rapport unique l'aurait révélé dès le premier
mois.
\end{exercise}

\section{Problème~: Trois chiffres, une journée}

\begin{problem}[{Problème du week-end --- réconciliation du P\&L d'un desk}]\label{pb:m1:pnl-and-positions:1}
Le nouveau quant du desk prend son poste. Le premier matin, trois chiffres de P\&L pour
la veille figurent sur la table du responsable de desk, qui en réclame un seul. Le desk
ne négocie qu'une action. Il a commencé la journée en \omterm{def:m1:pnl-and-positions:position}{position} longue de 12\,000 actions~; le
cours de clôture précédent était de 84.00. Exécutions de la veille, dans l'ordre~: vente de 7\,000 à 84.60~; vente de 9\,000 à
84.90~; achat de 6\,000 à 84.20~; achat de 10\,000 à 83.70. La dernière transaction de la journée
a été imprimée à 83.55, la cotation de clôture était de 83.70 en cours acheteur et 83.80 en cours vendeur, et le
prix de l'enchère de clôture était de 83.76. Les frais s'élevaient à 0.12 cent par action~; le financement et le coût de
l'emprunt ont coûté 310~\$.

\medskip\noindent\textbf{Partie I --- \omterm{def:m1:pnl-and-positions:position}{Position} et trésorerie.}
\begin{enumerate}
\item Donner la \omterm{def:m1:pnl-and-positions:position}{position} après chaque exécution et à la clôture.
\item Donner la trésorerie produite par les exécutions de la journée.
\item Calculer les frais.
\item Déterminer à quels moments le desk était vendeur, et de combien d'actions.
\end{enumerate}

\medskip\noindent\textbf{Partie II --- Trois valorisations.}
\begin{enumerate}[resume]
\item Donner le P\&L brut de la journée valorisé à la dernière transaction.
\item Puis au cours milieu de clôture.
\item Puis au prix de l'enchère de clôture.
\item Puis au prix qui serait obtenu en clôturant la \omterm{def:m1:pnl-and-positions:position}{position} contre
la cotation.
\item Quelle valorisation utiliser pour la prime du trader, et pourquoi~?
\end{enumerate}

\medskip\noindent\textbf{Partie III --- Explication} (utiliser le prix de l'enchère).
\begin{enumerate}[resume]
\item Calculer la ligne de la \omterm{def:m1:pnl-and-positions:position}{position} reportée.
\item Calculer la ligne des nouvelles transactions, exécution par exécution.
\item Ajouter les frais et le financement et donner le total net.
\item Vérifier que la part non expliquée est nulle.
\item Le trader affirme «~j'ai gagné de l'argent en négociant et je l'ai perdu sur la
\omterm{def:m1:pnl-and-positions:position}{position} overnight~». L'explication dit-elle cela~?
\end{enumerate}

\medskip\noindent\textbf{Partie IV --- Réalisé et latent.}
\begin{enumerate}[resume]
\item Les 12\,000 actions d'ouverture ont un \omterm{def:m1:pnl-and-positions:realised}{coût moyen} de 82.50. Suivre le
\omterm{def:m1:pnl-and-positions:realised}{coût moyen} et le \omterm{def:m1:pnl-and-positions:realised}{P\&L réalisé} à travers les quatre exécutions.
\item Donner le \omterm{def:m1:pnl-and-positions:realised}{P\&L latent} au prix de l'enchère.
\item Le réalisé plus le latent n'est pas le P\&L de la journée. Qu'est-ce que c'est, et que
faut-il soustraire pour retrouver le chiffre de la journée~?
\item Le relevé du courtier compensateur indique une \omterm{def:m1:pnl-and-positions:position}{position} de clôture de 13\,000
actions. Énumérer trois causes possibles, par ordre de vraisemblance, et ce qu'il faudrait
vérifier en premier.
\item Énoncer le \emph{résultat nommé}~: le P\&L net unique de la journée à remettre
au responsable de desk, et ses trois composantes les plus importantes.
\item En une phrase, énoncer la politique qui aurait empêché que trois chiffres
n'atteignent sa table.
\end{enumerate}
\end{problem}

\section{Questions d'entretien}

\begin{interviewq}[$\star$ \iqroles{trader, développeur, chercheur}]\label{iq:m1:pnl-and-positions:1}
Quelle est la différence entre le \omterm{def:m1:pnl-and-positions:realised}{P\&L réalisé} et le \omterm{def:m1:pnl-and-positions:realised}{P\&L latent}~? Le choix
entre premier entré, premier sorti et \omterm{def:m1:pnl-and-positions:realised}{coût moyen} change-t-il le montant d'argent que vous avez gagné~?
\end{interviewq}

\begin{interviewq}[$\star$ \iqroles{développeur}]\label{iq:m1:pnl-and-positions:2}
Pourquoi un système de P\&L ne devrait-il pas stocker les prix et la trésorerie sous forme de nombres à virgule flottante~?
Qu'utilisez-vous à la place~?
\end{interviewq}

\begin{interviewq}[$\star\star$ \iqroles{trader, chercheur}]\label{iq:m1:pnl-and-positions:3}
Vous êtes long de 10\,000 actions, la dernière transaction est à 50.30 et le marché est à 49.99
en cours acheteur, 50.01 en cours vendeur. Combien vaut votre \omterm{def:m1:pnl-and-positions:position}{position}~? Et si la \omterm{def:m1:pnl-and-positions:position}{position} comptait deux
millions d'actions~?
\end{interviewq}

\begin{interviewq}[$\star\star$ \iqroles{développeur, trader}]\label{iq:m1:pnl-and-positions:4}
Votre stratégie croit être à plat~; le fichier de fin de journée du courtier indique qu'elle est
longue de 400 actions. Décrivez comment vous déterminez lequel des deux a raison, et ce que le
système aurait dû faire dès l'instant où elles ont divergé.
\end{interviewq}

\begin{interviewq}[$\star\star$ \iqroles{chercheur, banque}]\label{iq:m1:pnl-and-positions:5}
Qu'est-ce qu'une explication du P\&L, et que vous apprend un montant non expliqué
persistant~?
\end{interviewq}

\begin{interviewq}[$\star\star\star$ \iqroles{développeur}]\label{iq:m1:pnl-and-positions:6}
Concevez le modèle de données d'un gestionnaire de \omterm{def:m1:pnl-and-positions:position}{positions} qui doit gérer des exécutions arrivant
dans le désordre, des transactions annulées et des corrections de prix. Qu'est-ce qui est
stocké, et qu'est-ce qui est dérivé~?
\end{interviewq}
